%% kapa: exam=1 round=19 year=2008 period=1 form=A subject=appraisal_law subject_ko=부동산관계법규 seq=27 number=67
%% source: Q-Net 감정평가사 기출문제 (cst00309) artlSeq=1001962 file=extracted/1008/1001962/2008년 제19회 감정평가사 1차시험 기출문제/1교시.hwp
%% conversion: kordoc 4.24.51 -> tools/gen1.py (auto); review=proofread-sonnet; flags=none
\begin{question}{67}
국유재산의 사용료에 관한 관련 법조문이다. 빈칸에 들어갈 내용으로 옳지 않은 것은?
\begin{framedtext}
\textbf{국유재산법 제25조(사용료) 제1항:} 행정재산 또는 보존재산의 사용·수익을 허가한 때에는 대통령령이 정하는 요율과 산출방법에 의하여 매년 사용료를 징수한다. \par
\textbf{국유재산법시행령 제26조(사용료율과 평가방법) 제1항:} 법 제25조제1항에 따른 연간 사용료는 시장 임대료를 고려하여 당해 재산의 가액에 다음 각호의 요율을 곱한 금액으로 하되, 월할 또는 일할계산할 수 있다. 다만, 다른 법령에서 국유재산의 사용료에 대하여 달리 규정하고 있는 경우에는 그 법령에 의한다. \par
1. 제24조 제1항 제1호(행정목적 또는 보존목적의 수행에 필요한 때)의 경우 : ( ① ) \par
2. 제24조 제1항 제2호(공무원의 후생목적을 위하여 필요한 때)의 경우 : ( ② ) \par
3. 경작용의 경우 : ( ③ ) \par
4. 주거용의 경우 : ( ④ ) \par
5. 기타의 경우 : ( ⑤ )
\end{framedtext}
\begin{choices}
  \item 1천분의 25 이상
  \item 1천분의 40 이상
  \item 1천분의 25 이상
  \item 1천분의 25 이상
  \item 1천분의 50 이상
\end{choices}
\end{question}
\answer{3}
